\documentclass[../../main.tex]{subfiles}% 注意这里的文件路径不能用 ./main.tex ,否则用latexmk编译子文件会报错
\graphicspath{{\subfix{./image/}}} % 指定图片目录,后续可以直接使用图片文件名
% 注意这里的文件路径不能用 ../../image/ ,否则用latexmk编译子文件会报错

% 例如:
% \begin{figure}[H]
% \centering
% \includegraphics[scale=0.3]{图.png}
% \caption{}
% \label{figure:图}
% \end{figure}
% 注意:上述\label{}一定要放在\caption{}之后,否则引用图片序号会只会显示??.

\begin{document}

\subsection{自由群和群的表现}

\begin{example}
若 $n \geqslant 3$, 试问 $A_n \times \mathbb{Z}_2$ 与 $S_n$ 是否同构?
\end{example}
\begin{solution}
注意到$\{1\}\times \mathbb{Z}_2$是$A_n \times \mathbb{Z}_2$的2阶正规子群.
由\refpro{proposition:对称群正规子群的分类}知,当 $n \geqslant 3$ 时, $S_n$ 无 2 阶正规子群. 故 $A_n \times \mathbb{Z}_2$ 与 $S_n$ 当然不同构.

\end{solution}

\begin{example}
设 $G = G_1 \times \cdots \times G_n$, $H$ 是 $G$ 的子群. 问 $H$ 是否一定形如 $H = H_1 \times \cdots \times H_n$, 其中 $H_i \leqslant G_i$?
\end{example}
\begin{solution}
否.例如, $K_4 = \{(1),(12)(34),(13)(24),(14)(23)\} = \{(1),(12)(34)\} \times \{(1),(13)(24)\}$(同构意义下). 显然$\{(1),(14)(23)\}$ 是 $K_4$ 的子群,但$\{(1),(14)(23)\}$无法写成$H_1\times H_2$,其中$H_1\leqslant \{(1),(12)(34)\},H_2\leqslant \{(1),(13)(24)\}$.

否则,设$\{(1),(14)(23)\}=H_1\times H_2$,其中$H_1\leqslant \{(1),(12)(34)\},H_2\leqslant \{(1),(13)(24)\}$.

(i)若$|H_1|$或$|H_2|$为1,则$H_1\times H_2 = \{(1),(12)(34)\}$或$\{(1),(13)(24)\}\neq \{(1),(14)(23)\}$,矛盾!

(ii)若$|H_1|=|H_2|=2$,则$|H_1\times H_2|=4$,但$|\{(1),(14)(23)\}|=2$,矛盾!

\end{solution}

\begin{example}
设 $G = G_1 \times G_2$, $H \lhd G$ 且 $H \cap G_i = \{1\}$, $i=1,2$. 试证 $H \leqslant Z(G)$. 特别地, $H$ 是 Abel 群.
\end{example}
\begin{proof}
因为 $H \lhd G$, $G_i \lhd G$, $H \cap G_i = \{1\}$, 故由\rreflem{lemma:换位子群的基本性质}{lemma:换位子群的基本性质-3}知$H$ 中任一元与 $G_i$ 中任一元可换, 因此 $H$ 中任一元与 $G$ 中任一元可换, 即 $H \leqslant Z(G)$.

\end{proof}

\begin{example}
令$G = G_1 \times \cdots \times G_n$, 且对任意$i \neq j, 1 \leqslant i,j \leqslant n, |G_i|$和$|G_j|$互素。则$G$的任意子群$H$都是它的子群$H \cap G_i\ (i = 1, \cdots, n)$的直积。
\end{example}
\begin{proof}
设$h \in H$, $h = h_1 \cdots h_n$, $h_i \in G_i, i = 1,2,\cdots,n$。由\rrefthe{theorem:群的内直积分解的基本性质}{theorem:群的内直积分解的基本性质-1}知$G_i$中任一元与$G_j(i \neq j)$中任一元可换，故对任一整数$t$, 有$h^t = h_1^t \cdots h_n^t$。令$a = |G_i|$, $b = \prod\limits_{\substack{j \neq i}} |G_j|$, 则$h_i^a = 1$, $h_j^b = 1\ (j \neq i)$。由条件知$(a,b) = 1$, 故有整数$l$和$k$, 使得$la+kb = 1$。故
\begin{align*}
h_i = h_i^{la+kb} = h_i^{kb} = h^{kb} \in H \cap G_i, \quad i = 1, \cdots, n.
\end{align*}
从而
\begin{align*}
H = (H \cap G_1) \cdots (H \cap G_n) = (H \cap G_1) \times \cdots \times (H \cap G_n).
\end{align*}

\end{proof}





















\end{document}